\chapter{国内外研究现状}

组合设备（cluster tool，CT）调度研究起初关心连续加工同一类晶圆时，设备每隔多久能够产出一片晶圆。随着同一批次内出现多种工艺路线、一个腔室容纳多片晶圆，或者加工中途需要清洁腔室，决定产出的资源也会随设备状态而改变。对本文研究的旋转腔设备，有限批次的交付还取决于产品片、补位用的辅助片和清洁操作如何共用机械手，以及最后一片产品何时回到装载端口。下面沿着设备调度问题和求解方法两条线索梳理国内外研究。

\section{稳态吞吐率与驻留约束下的可调度性}

早期吞吐率模型把晶圆在腔室中接受工艺处理的时间与机械手搬运晶圆的时间放到同一周期内考察。1994年，Perkinson等分别计算单片和整批晶圆的加工时间，分析了工艺处理时间、搬运时间与设备最大吞吐率之间的关系\citep{perkinson1994throughput}。在机械手有充分时间服务各腔室时，缩短瓶颈工序的处理时间可以提高产出。当机械手已连续执行取放和移动动作时，即使增加腔室处理能力，晶圆仍需等候搬运。这使设备能力评价从单个腔室的标称加工时间转向腔室与机械手共同形成的周期。

晶圆加工结束后还可能受到残余气体或温度的影响，设备需要在工艺允许的时间窗内将其取出。于是，达到给定吞吐率的动作序列还要通过可调度性检验，即能否为该序列安排一组同时满足资源互斥和驻留上限的动作时刻。2002年，Rostami和Hamidzadeh同时考虑工艺模块与搬运模块的驻留限制，研究时间窗对设备调度的影响\citep{rostami2002optimal}。2010年，伍乃骐和周孟初针对双臂组合设备的稳态运行，从机械手等待时间推导可调度条件及最优稳态调度的闭式解\citep{wu2010closedform}。由此可以先判断重复动作序列能否在时限内执行，再计算可行序列所能达到的周期长度。

晶圆重入同一模块时，当前加工还要为同片晶圆后续再次访问该模块留下时间。2014年，乔岩等研究了原子层沉积中带重入和驻留约束的双臂组合设备，分析重复访问对机械手调度的影响\citep{qiao2014revisiting}。次年，乔岩在广东工业大学完成了关于半导体组合设备建模、优化调度与控制的博士论文\citep{qiao2015thesis}。对于更一般的离散事件系统，周家忠在西安电子科技大学的博士论文中采用时间扩展可达图表示操作顺序和时刻，通过集束搜索或状态聚类缩小搜索范围\citep{zhou2022thesis}。前一类工作利用特定设备结构推导时序性质，后一类方法则为动作顺序尚待搜索的情形提供了状态空间表示。袁凤连等在《自动化学报》的综述中归纳了这条研究脉络，并将驻留限制、作业时间波动、重入、多品种、模块故障、清洗和组合设备群列为七类典型问题\citep{yuan2023survey}。

\section{多品种投放、有限批次与多工位腔室}

从单一配方扩展到多品种后，即使各品种单独运行时都满足驻留上限，混合投放后仍要重新安排机械手。王际鹏等在2020年针对双臂组合设备提出多路交换序列，给出多种晶圆流模式下周期运行的可调度条件\citep{wang2020multiple}。当设备同时接收两类晶圆时，投放顺序还会改变并行腔室的负载。即使两种投放方案在一个周期内送入相同数量的晶圆，长工时晶圆集中到同一腔室，也会使另一腔室先完成加工并等待机械手。韩国研究者Ko等利用赋时事件图分析双臂设备，在保持最小周期时间的条件下，通过调整多类晶圆的投放顺序均衡并行腔室的负载，减少加工后的腔内等待\citep{ko2021balance}。2025年，Yang等将两类晶圆的投放顺序和单臂机械手的任务顺序同时作为混合整数规划的决策，分别评价周期吞吐率与驻留时间约束满足率\citep{yang2025miptwowafers}。因此，投放次序既决定哪一种产品先接受加工，也决定机械手随后需要优先服务哪一个腔室。

小批量生产中，设备尚未进入稳态，一个批次便可能接近结束。周期模型省去的首片入腔和末片离腔过程，此时会占据相当一部分批次时间。李金成和潘春荣研究单臂设备插入紧急订单后的暂态过程，针对不同加工参数推导插单可行条件及机械手等待时间\citep{li2023urgent}。罗亚波等进一步研究双臂设备的终止暂态，利用仍在加工末批晶圆的时间完成部分换模操作，在满足驻留约束的前提下缩短了算例中的换批时间\citep{luo2026mask}。从完整批次出发，Yang等也已针对两类晶圆建立单臂设备非周期调度模型，覆盖启动至收尾过程\citep{yang2025noncyclic}。因此，描述有限批次时，工艺路线之外还应明确完工事件取在末片离腔、完成后续搬运，还是回到原装载端口。

设备结构的扩展沿着不同的方向发展。多个并行腔室分别接收晶圆时，主要决策是如何分配负载。两个组合设备通过缓冲模块连接时，还要协调两台机械手对交接位置的访问。Kim在2023年的博士论文中用赋时事件图分析串联组合设备，再将反馈弧引入多周期序列的分析，以识别中间缓冲位置的阻塞及其造成的晶圆延迟\citep{kim2023tandemphd}。潘春荣等针对带全并行加工模块的双组合设备，利用面向资源的Petri网分析机械手任务冲突，给出了相应的可调度条件和稳态周期求解方法\citep{pan2024twin}。

将多个工位设在同一工艺腔内，又会引入不同的取片问题。Jee在韩国科学技术院完成的博士论文研究多槽位设备，考察一次处理多片晶圆后各片等待机械手取出的时间，并将多臂机械手的动作次序和含腔室清洁的设备配置纳入分析\citep{jee2021thesis}。Zou等研究有驻留时间约束的双工位单臂组合设备，进一步将同一模块内的多个工位纳入调度\citep{zou2024twospace}。

前端供片也会改变真空侧的可行时序。Huang等研究由设备前端、装载锁和真空模块组成的单臂设备，将前端搬运与真空侧时序放在同一设备内求解，并要求晶圆满足工艺驻留上限\citep{huang2025efem}。本文研究的双源分别供应产品晶圆和辅助晶圆。双源配置下，产品片从大气侧经装载锁进入真空侧，辅助片从真空存储区取用。产品品种决定加工需求，晶圆来源决定它占用的运输路径。两片晶圆即使能参加同一次加工，也必须能在同一装片窗口到达旋转腔的可访问工位。

\section{腔室清洁与辅助晶圆的复用}

安排清洁时，首先要区分清洁何时触发，以及触发后其他腔室如何继续生产。Yu和Lee研究双臂设备的腔室清洁，分析清洁周期内的晶圆延迟，通过减少向并行腔室装入的片数和调整机械手动作开始时刻控制等待\citep{yu2020delay}。乔岩等在双臂设备的周期序列中插入虚拟晶圆，使相应腔室能在该时段进行清洁，再用二元整数规划优化吞吐率\citep{qiao2022binary}。镇璐等在《系统工程理论与实践》讨论清洁时间较长的情形，为单臂设备设计两种机械手调度策略，并按工序分配等待时间以满足驻留限制\citep{zhen2023cnclean}。

对于按腔室状态触发的清洁，Li等研究了单臂设备的调度，Zhen等进一步处理双臂设备，给出在不同清洁情形下满足驻留约束并恢复稳态运行的策略\citep{li2022conditioncleaning,zhen2025cleaning}。多品种与清洁也已有联合研究，Yang等在两类晶圆、驻留限制及清洁并存的单臂设备上联合决定投放与机械手任务次序\citep{yang2025twoclean}。能否利用一段空闲时间清洁腔室，取决于清洁持续多久，也取决于这段时间内其他产品是否已到必须卸片的时刻。

在这些时序模型中，虚拟晶圆可以用来保留机械手动作序列中的位置。例如，罗亚波等用虚拟晶圆描述终止暂态中的换模时机\citep{luo2026mask}。制造现场的控片和挡片则有实物库存。Chen和Lee研究晶圆厂控片与挡片的补充、降级和回用，说明循环使用的非产品片同样需要库存控制\citep{chen2003dummy}。杨懿淑在新竹交通大学完成的博士论文进一步考虑多期需求的不确定性，通过随机规划确定新片供给和降级数量，并用服务水平约束控制缺片风险\citep{yang2009dummyphd}。这些工作回答了晶圆厂应准备多少片、怎样补充的问题。在设备调度中，还需知道已经借出的某一片何时能再次使用。

在本文的满片加工条件下，产品片数量不足时由辅助片占据空余工位，清洁时全部工位都要装入辅助片。于是，生产末端的补位会影响下一次清洁最早何时可以装片。即使两座腔室均已空闲，如果可用辅助片只够其中一座执行满片清洁，两次装片也要错开。描述这一约束需要把每片辅助片的源端取片、腔内占用和归还时刻连接起来。晶圆厂级的库存数量给出总量边界，设备动作序列给出每一时刻实际可借用的数量。

\section{驻留时间、路线等待与批内节拍}

调度方案的时间特性取决于测量的起止事件。缩短晶圆从加工结束到卸片的间隔，直接服务于驻留要求。在单周期产量相同的条件下，缩短平均开工间隔可以提高产出率，而缩小相邻间隔的差异则能使连续生产中的供片与取片更有规律。这几个指标可以在同一调度方案中呈现不同变化。朱清华等在双臂组合设备群上研究加工后停留时间，在取得最高吞吐率的同时优化工序加工后的总停留时间，并进一步用线性规划平衡相邻设备的停留时间差异\citep{zhu2023regulation}。2026年，Kim和Yu进一步以单片晶圆流经整台设备的时间为指标，研究调节设备内在制晶圆数的效果。在他们分析的设备中，降低装载水平可缩短单片流动时间，同时可能延长设备周期\citep{kim2026flow}。这两类结果说明，逐片时间可以在吞吐率相同的条件下得到改善，也可能需要以一部分产出为代价，具体取决于可以调整哪一类决策。

该区分在动作时间发生波动时尤其重要。Kim等研究随机时间扰动下的设备反馈控制，分析了设备受扰后晶圆停留时间能够重新稳定的条件\citep{kim2022feedback}。潘春荣等则在给定加工时间概率分布的条件下研究带并行腔室和周期清洁的设备，给出单臂设备平均周期时间的计算方法和双臂设备相应指标的近似算法\citep{pan2024variation}。前者处理执行中已经发生的偏差，后者预测许多周期的平均表现。本文给定各动作的标称时长，考察同一有限批次中由资源竞争造成的等待和间隔变化。

具体而言，本文以最后一片产品回到原装载端口的时刻衡量批次完工。对已经取得的可行排程，先保留其完工时间上界和主要离散决策，再调整连续时刻与允许变化的取放同步选择。路线等待扣除必要搬运或操作时长，分别统计总量和最大单项。链内节拍比较同一腔室、同一条$2\times2$加工链内连续三个生产周期形成的前后两个开始间隔。这样，同一份排程即可分别说明交付时刻、途中等待和链内节拍，清洁前后的间隔也不会混入同一条链的节拍评价。第三章式~\eqref{eq:primary-objective-cn}和式~\eqref{eq:secondary-rhythm-cn}给出了相应的计量边界。

\section{设备动作学习与混合整数规划选割}

组合设备中的强化学习首先可以作用于机械手。Kim和Lee研究两种晶圆在单臂设备上的并发加工，以强化学习和自适应搜索同时安排投放次序与机械手任务，使策略可以在预设周期序列之外选择动作\citep{kim2025adaptivecluster}。Kim和Kim针对非周期组合设备调度，将晶圆配方信息用于动作决策，结合当前调度状态选择下一步机械手动作以避免死锁\citep{kim2025noncyclicdeadlock}。2026年，Lai等进一步研究带设备前端的单臂设备，在由装载锁连接的大气侧和真空侧上开展非周期调度\citep{lai2026front}。在这些方法中，策略的一次选择会改变下一步哪些晶圆可以移动、哪些资源可供使用。

另一条路线先将设备约束写成混合整数线性规划（MILP），在分支定界搜索中用学习方法选择割平面。有效割可以排除线性松弛中的部分分数解，同时保留原整数可行解。Huang等将候选割选择表述为排序问题，利用多实例学习从候选割特征中估计优先级\citep{huang2022cutranking}。Paulus等以前瞻计算得到的界改善作为监督信号，用模仿学习近似开销较高的选割准则\citep{paulus2022looking}。两者处理的都是求解器已经生成的候选割，评价所选割时需要考虑松弛界改善及其计算代价。

Wang等在2023年的层次序列模型（HEM）中用上层策略确定选割数量，再用下层策略依次选出相应数量的割。2024年的期刊扩展研究了序列和集合两种候选表示，以及已选割之间的相互影响\citep{wang2023hem,wang2024hemset}。这一进展与旋转腔调度的联系在于，两个单独评分较高的割可能都集中在同一组腔室分配变量上，而候选集中还有涉及库存或机械手先后次序的割。第四章在HEM基础上提取变量角色特征并重排候选割。此处策略的输出是求解器采用的割集合与顺序，逐片排程仍从满足设备约束的整数可行解中解码。

\section{可行排程构造与原始解的邻域改进}

对最小化产品完工时间的模型，每一份完整可行排程都给出一个上界。寻找首个可行排程与改进已有排程可以利用不同的信息。前者可以从最难安排的资源入手。周炳海和黎明研究带并行腔的双集束型设备，围绕产能约束资源组织启发式调度，在多品种和驻留限制下协调并行腔与机械手\citep{zhou2016ccr}。对后者，已有排程还指明了具体的资源使用次序。例如，最后回片的产品若在装载锁前等待，缩短另一腔室一项与该产品无资源关联的动作，对其回片时刻没有影响。要使该产品提前回片，可能需要挪动此前占用同一装载锁或机械手的另一片晶圆。

通用整数规划为保留可行排程的大部分结构提供了两类邻域。2003年，Fischetti和Lodi提出局部分支，在当前整数解附近限制二元变量允许翻转的数量，由求解器在该范围内寻找改进解\citep{fischetti2003local}。Danna等提出松弛诱导邻域搜索（RINS），固定当前整数可行解与线性松弛解中取值一致的整数变量，再求解剩余子问题\citep{danna2005rins}。局部分支先限定改动幅度，RINS则依据两组解的一致性确定保留哪些决策。对具有实体库存的排程，移动一个产品对到另一腔室可能同时改变清洁触发位置、辅助片归还时刻和机械手次序。若邻域只允许改变分腔变量，相关动作就可能没有重新配合的余地。

邻域大小会影响子问题的求解成本。Hendel在SCIP中根据各大邻域启发式以往的收益选择下一个邻域，并通过固定变量调节子问题规模\citep{hendel2022alns}。Liu在蒙特利尔综合理工学院的博士论文中用学习方法确定局部分支的邻域大小，并利用优质解的结构特征选择变量生成邻域\citep{liu2022thesis}。对旋转腔设备，可进一步沿产品回片路径查找造成延迟的资源，联合放开有关的分配、占位与搬运次序，再用原MILP约束修复排程。可按相同计算预算下的可行上界和子问题耗时评价该邻域。本文的结构化初始化从配对、分腔和周期安排出发，为主问题建立初始上界。

\section{研究评述}

从稳态周期分析到多品种非周期调度，国内外研究已能在多种设备配置下联合确定投片和机械手动作，处理清洁、驻留及前端供片约束。控片与挡片的库存研究也已说明，可回用晶圆的供给会约束生产。对于本文的四工位旋转腔，这些问题在同一次换片中相遇。采用$4\times1$并行加工时，四个工位同批处理。采用$2\times2$链式加工时，一对晶圆要留在腔内继续处理，另一对要在旋转后换入或换出。模式选择因而同时决定哪一面可装卸、需要多少辅助片，以及两腔何时争用同一机械手。仅根据单腔的稳态周期或库存总量，无法还原这一动作序列。

表~\ref{tab:review-decisions}比较了与本文最接近的几项研究。是否将辅助片作为可返回的实体、是否记录每次取放，对问题的决策范围都有直接影响。

\begin{table}[htbp]
\centering
\caption{代表性研究的决策范围与评价对象}\label{tab:review-decisions}
\small
\begin{tabularx}{\textwidth}{@{}p{0.22\textwidth}XX@{}}
\toprule
研究 & 主要决策 & 评价对象 \\
\midrule
杨懿淑\citep{yang2009dummyphd} & 各期新片供给量、各等级控片与挡片的降级数量 & 不确定需求下的库存成本与服务水平 \\
Jee\citep{jee2021thesis} & 多槽位腔室的取放次序、机械手与腔室配置 & 同腔多片晶圆加工后的取片延迟 \\
镇璐等\citep{zhen2023cnclean} & 有清洁作业时各工序的机械手等待时间 & 驻留时限内的可行调度 \\
Yang等\citep{yang2025noncyclic} & 启动至收尾期间两类晶圆的投放和搬运顺序 & 有限批次的完成时间 \\
本文第三章 & 晶圆配对、分腔、旋转加工、辅助片借还与共享资源次序 & 最后产品回片时刻、路线等待与链内节拍 \\
\bottomrule
\end{tabularx}
\end{table}

第三章以晶圆对为单位连接周期占位、辅助片借还与逐片供回路线，以产品回片界定批次完工。第四章研究如何构造初始排程，以及如何通过选割改进后续搜索。在相同计算预算下比较初始上界、最终上界和求解下界，可分别评价初始排程的质量、排程的后续改进和下界的收紧程度。
